\documentclass[11pt]{article}

\usepackage[margin=1in]{geometry}
\usepackage{amsmath}
\usepackage{booktabs}
\usepackage[colorlinks=true,linkcolor=blue,urlcolor=blue]{hyperref}

\setlength{\emergencystretch}{2em}
\newcommand{\ngvla}{ngVLA}
\newcommand{\arcsec}{^{\prime\prime}}

\title{How the \ngvla\ Subarray Calculator Constructs a Ranked List}
\author{Short technical note accompanying the subarray-selection memo}
\date{21 August 2026}

\begin{document}
\maketitle

\begin{abstract}
This note describes the query-time construction of the ranked subarray list,
assuming that the robust--taper measurements and subarray metadata already
exist.  A 30~GHz request is used so that the frequency scaling is transparent.
The main point is that the calculator does not rank subarrays by inefficiency
$\eta$ alone.  It first finds one feasible imaging operating point for each
subarray, then ranks those points by integration time or antenna hours.
\end{abstract}

\section{Inputs and conventions}

For each offered subarray $i$, the calculator needs metadata
$(N_i,b_{\min,i},C_{{\rm LAS},i},\theta_{{\rm cap},i})$ and measured lattice
points
\[
  p=(\theta,\eta,s,{\rm weighting},{\rm robust},{\rm taper}).
\]
Here $s$ is the PSF skirt at one fitted FWHM, and
\begin{equation}
  \eta(p)=\frac{\sigma_{\rm dirty}(p)}
  {\sigma_{\rm dirty}({\rm natural},0)}
\end{equation}
is normalized to the natural, untapered rms of the \emph{same} subarray.
Thus equal $\eta$ does not imply equal absolute sensitivity for subarrays with
different antenna counts.

The registered configuration and observing mode select one metadata table and
one curve table.  The default is Rev~F, four-hour track.  The selector evaluates
every metadata row independently; it does not first construct a single envelope
across different subarrays.

At a general observing frequency $f$, the reference-curve coordinate is
\[
  \theta_{30}=\theta_{\rm req}\frac{f}{30~{\rm GHz}}.
\]
For the 30~GHz case in this note, $\theta_{30}=\theta_{\rm req}$ and returned
tapers also need no scaling.  The current model treats $\eta$ and skirt as
frequency-invariant.

\section{One candidate from each subarray}

Let $m$ be the fractional beam tolerance (currently 0.05 by default).  The
allowed reference-frequency beam interval is
\begin{equation}
  I_\theta=[\theta_{\rm req}(1-m),\theta_{\rm req}(1+m)].
\end{equation}
The achieved beam can therefore lie at either edge of the interval; it need not
equal the central request.

The measured robust--taper lattice is made continuous conservatively.  Adjacent
points are interpolated along fixed-robust taper rows and fixed-taper Briggs
columns.  Valid four-corner Briggs cells also admit interior robust and taper
values.  On an edge, with
\[
 u=\frac{\ln(\theta/\theta_a)}{\ln(\theta_b/\theta_a)},
\]
the engine interpolates $\ln\eta$ linearly, interpolates skirt linearly, and
normally interpolates taper in log space.  It does not extrapolate beyond a
measured family or bridge a pruned gap.

For a point-source request, the subarray's operating point is
\begin{equation}
 p_i^*=\mathop{\rm arg\,min}_{p:\,\theta(p)\in I_\theta}\eta(p).
\end{equation}
For an imaging request with required skirt $s_0$, the current per-subarray
policy is:
\begin{enumerate}
  \item If any point in $I_\theta$ has $s\le s_0$, choose the lowest-$\eta$
        such point.  This is a \emph{strict} result.
  \item Otherwise, if relaxation is enabled, choose the best achievable skirt
        in $I_\theta$, breaking skirt ties by $\eta$.  This is a
        \emph{relaxed} result and does not meet the requested skirt.
  \item With relaxation disabled, absence of a strict point excludes the
        subarray.
\end{enumerate}
Interpolation provenance and the winning \texttt{tclean} weighting, robust,
and taper are retained in the returned candidate.

Two other feasibility gates apply.  A heterogeneous subarray is excluded when
the entire beam window is coarser than its beam cap; if the interval only
overlaps the cap, it is clipped at the cap.  A requested largest angular scale
is checked using
\begin{equation}
 {\rm LAS}_i=C_{{\rm LAS},i}\frac{\lambda}{b_{\min,i}}206264.8.
\end{equation}
At 30~GHz and $C_{\rm LAS}=0.8$, this is approximately
$1649/b_{\min,i}[{\rm m}]$ arcsec.  LAS is reported but is not a gate when the
request omits it.  There is no query-time $\eta\le5$ rejection.

\section{Sensitivity cost and the role of antenna count}

For a supplied flux-density goal $g$, the sensitivity wrapper obtains the
one-hour natural rms $\sigma_{{\rm nat},i}$ from the Memo~96 calculator and
applies the curve-derived inefficiency externally:
\begin{equation}
 \sigma_i(t)=\eta_i\sigma_{{\rm nat},i}(1~{\rm h})
             \sqrt{\frac{1~{\rm h}}{t}},
 \qquad
 t_i=\left(\frac{\eta_i\sigma_{{\rm nat},i}(1~{\rm h})}{g}\right)^2{\rm h}.
\end{equation}
The resource cost is then
\begin{equation}
 A_i=N_i t_i
\end{equation}
in antenna hours.  A brightness-temperature goal is converted to Jy/beam
using the candidate's achieved beam and request frequency before evaluating
the same expression.

There is no double counting of $N_i$ in these equations.  The upstream natural
rms already contains
\begin{equation}
 \sigma_{{\rm nat},i}\mathrel{\propto}
 \frac{1}{\sqrt{N_i(N_i-1)t}}.
\end{equation}
The wrapper rescales this reference natural sensitivity only when the backing
calculator does not already represent the candidate being evaluated.  There
are two intelligible cases.  First, a named backing entry can carry the antenna
count of an older configuration: for example, the calculator has a
Rev-D-vintage core entry, not a separate Rev~C core entry, so a Rev~C request
must replace the stored count by the Rev~C metadata count.  Second, some
offered subarrays have no named backing entry at all; for these, the
main-array natural sensitivity is used as a same-dish reference and rescaled
to the candidate's antenna count.  Ordinary Rev~F and Rev~D named entries have
matching stored and actual counts, so their correction factor is exactly one.
If both a donor reference and a revision-specific count are involved, the
intermediate count cancels algebraically, leaving one main-to-actual
correction.  Finally, multiplying $t_i$ by $N_i$ converts elapsed time to the
EOP resource unit; it does not alter the thermal noise.

When no numerical sensitivity goal is supplied, the calculator defines the
common fiducial goal
\begin{equation}
 g_{\rm fid}=\sigma_{{\rm nat},{\sc main}}(1~{\rm minute})
\end{equation}
at the requested frequency and spectral setup, and evaluates every subarray
with the same absolute time equation above.  Thus natural, untapered MAIN
takes one minute by construction, while other beams and subarrays have finite,
interpretable costs.  This is still not simply a lowest-$\eta$ ordering: the
natural-rms anchor includes $N_i(N_i-1)$, so a larger subarray can rank ahead
despite a larger taper penalty.

\section{Final ordering and current caveats}

For imaging, feasible candidates are first partitioned into strict results
that meet $s_0$ and relaxed results that do not.  Strict results are always
listed first.  Within each group, the default sort is increasing antenna
hours, with fewer antennas as a tie breaker; the alternative is increasing
integration time.  If no feasible subarray meets $s_0$, the calculator emits
a structured warning and retains the relaxed list.  Python callers can
explicitly request a mixed-pool policy when a custom ranking needs to trade
skirt compliance against cost.

Sensitivity failures cannot enter a relative-cost fallback.  Failure to cost
an explicit science goal excludes the affected candidate.  Failure of the
implicit fiducial leaves the independently determined beam/LAS feasibility
visible, but disables the entire recommendation ranking and returns a
structured cost warning.  Excluded subarrays follow the feasible rows with
their reasons.

\section*{Implementation locations}

The relevant implementation is in \path{sa_select/curves.py},
\path{sa_select/envelope.py}, \path{sa_select/selector.py},
\path{sa_select/heuristics.py}, and \path{sa_select/sensitivity.py}.  The
browser implementation is in \path{curvegen/report_template.html}, with its
data payload assembled by \path{curvegen/make_report.py}.

\end{document}
